\section{Conclusiones}
\label{sec:conclusiones}

\subsection{Sobre los objetivos}

Los ocho objetivos específicos de la sección~2 se cumplieron, y siete de los
ocho tienen verificación automatizada asociada. Las seis pruebas funcionales
exigidas pasan, y todas comprueban el resultado en la base de datos en lugar
de en las respuestas HTTP.

\subsection{Sobre las garantías}

La conclusión técnica central es que \textbf{las garantías de integridad no
dependieron de la estrategia de concurrencia elegida}. En las cuatro carreras
de escritura, con cien usuarios durante sesenta segundos y cuatro procesos
compitiendo, el libro mayor sumó exactamente cero, ninguna cuenta quedó
descuadrada y ningún saldo resultó negativo, tanto con bloqueo pesimista como
con \codigo{SERIALIZABLE}.

Lo que sí dependió de la estrategia fue cuánto trabajo se completó: 148,3
frente a 10,4 transferencias confirmadas por segundo bajo contención máxima.
Corrección y rendimiento resultaron ser dimensiones separables, y conviene no
confundirlas al justificar una decisión.

\subsection{Sobre dónde deben vivir las invariantes}

Nueve de las diez invariantes del sistema se aplican en PostgreSQL y no en
Python. Las diecisiete pruebas de esquema las atacan con SQL directo, sin
pasar por la aplicación: si alguna viviera solo en el código, esas pruebas se
pondrían rojas.

Esa decisión tuvo un rendimiento concreto. La restricción de no negatividad
nunca llegó a dispararse durante la actualización perdida, porque el saldo
calculado a partir del valor obsoleto no cruzaba el cero. Una restricción de
motor acota el daño, pero no sustituye a la prueba. Ambas cosas eran
necesarias.

\subsection{Sobre el método}

Los dos defectos más valiosos del proyecto no los encontró la revisión de
código, sino las pruebas que ejercitan el sistema en las condiciones en las
que falla:

\begin{itemize}[leftmargin=1.2em]
  \item La \textbf{prueba de concurrencia} destapó una actualización perdida
    con la que las 76 pruebas existentes pasaban todas. Tomar el bloqueo
    correcto en SQL no basta si la capa de mapeo devuelve datos cacheados.
  \item La \textbf{prueba de carga} destapó un agotamiento del pool de
    conexiones causado por retener la sesión fuera del hilo que la usaba.
    Corregirlo multiplicó por 7,4 el rendimiento de lectura.
\end{itemize}

A eso se añade una tercera lección, menos cómoda: una de las mediciones de
este informe hubo que descartarla por no ser válida. Revisar críticamente los
propios datos antes de publicarlos resultó ser tan necesario como medirlos.

\subsection{Sobre lo que no se hizo}

PayCore no incorpora Redis, ni Kafka, ni Kubernetes, ni microservicios, ni
modelos de aprendizaje automático. Para cada uno está escrita la condición
concreta bajo la cual entraría. Esa lista de exclusiones, con su justificación
y su condición de reingreso, es tan parte del diseño como lo que sí se
construyó: una decisión técnica no se defiende solo explicando lo que se
eligió, sino también lo que se descartó y por qué.

\subsection{Valoración final}

El sistema cumple su propósito: demostrar, con pruebas reproducibles y
mediciones propias, que las garantías que un motor de pagos debe ofrecer se
pueden construir y verificar. No es una plataforma bancaria de producción y el
informe dedica una sección entera a explicar por qué no lo es. Entre un
prototipo que exagera lo que hace y uno que delimita con precisión lo que
garantiza, el segundo es el que se puede defender.
